% Generated draft: requires mathematical and source review.
\documentclass[11pt]{article}
\usepackage[margin=25mm]{geometry}
\usepackage{fontspec}
\setmainfont{Latin Modern Roman}
\usepackage{amsmath,amssymb,mathtools,unicode-math}
\setmathfont{Latin Modern Math}
\usepackage{enumitem,needspace,xurl,hyperref}
\setlength{\parindent}{0pt}
\setlength{\parskip}{5pt}
\setlength{\emergencystretch}{3em}
\newcommand{\N}{\mathbb{N}}
\newcommand{\Z}{\mathbb{Z}}
\newcommand{\Q}{\mathbb{Q}}
\newcommand{\R}{\mathbb{R}}
\newcommand{\C}{\mathbb{C}}
\newcommand{\F}{\mathbb{F}}
\newcommand{\A}{\mathbb{A}}
\newcommand{\PP}{\mathbb{P}}
\newcommand{\E}{\mathbb{E}}
\newcommand{\eps}{\varepsilon}
\newcommand{\dd}{\,\mathrm{d}}
\begin{document}
\section*{UnsolvedMath ID: 30004338}
\textbf{Source problem code:} catalogue entry
\subsection{Definitions and mathematical statement}
Let \((X,\mu)\) be a standard probability space, and let \(\operatorname{Aut}(X,\mu)\) denote its group of invertible measure-preserving transformations, identified modulo equality almost everywhere. The source does not assume that \((X,\mu)\) is nonatomic. It presents the following two separate questions.

\begin{enumerate}
\item How many essentially free quasi-actions of \(\operatorname{Aut}(X,\mu)\) are there up to conjugacy? Does every such quasi-action arise as a factor of the canonical product quasi-action denoted by \(\rho_\infty\)?

The supplied record does not define a quasi-action, its target space, the applicable meaning of essential freeness or conjugacy, the factor relation, or the construction of \(\rho_\infty\). Consequently, the objects being classified and the proposed universal factor property cannot yet be specified precisely. These terms require their source-specific definitions; no definition of quasi-action is selected here.

\item For an integer \(k\geq 2\), write
\[
 r_k=\frac{\sqrt{2k-1}}{k}.
\]
The source asks whether, for every such \(k\) and every \(\varepsilon>0\), there exists \(C\) such that every \(k\)-generated subgroup \(H\) of a free group with ``no nontrivial reduced-word relations'' satisfies
\[
 \lVert A_H\rVert\leq r_k+\varepsilon,
\]
where \(A_H\) is the normalized adjacency operator referred to in the source.

For clarity about the relation terminology, relative to a chosen generating tuple \((s_1,\ldots,s_k)\), a reduced word is a word in formal letters \(a_1^{\pm1},\ldots,a_k^{\pm1}\) with no adjacent inverse pair. A nontrivial reduced-word relation is a nonempty such word whose evaluation at \((s_1,\ldots,s_k)\) is the identity. This terminology alone does not specify which generating tuples or word lengths the source intends to quantify over.

The domain and role of \(C\) are absent: it occurs in neither the stated hypothesis nor the inequality. In particular, the record gives no word-length or girth threshold involving \(C\). Its background describes a large-girth Cayley-graph question, whereas its displayed question excludes nontrivial reduced-word relations without any length restriction. These formulations cannot silently be identified. The generating set defining the Cayley graph, the normalization of its adjacency operator, and the space on which the operator norm is taken are also unspecified. The displayed inequality therefore records the requested bound without supplying a complete operator-theoretic formulation.
\end{enumerate}

Both questions are retained under the single supplied permanent ID \(30004338\). This is a statement-only draft with unresolved formulation issues, requiring human mathematical review.
\subsection{Short English statement}
The record asks for a classification of essentially free quasi-actions of \(\operatorname{Aut}(X,\mu)\), including whether all are factors of \(\rho_\infty\), and separately asks for an almost-Ramanujan adjacency bound for certain subgroups of free groups. Definitions needed for the first question are missing. The second question leaves \(C\) unused and does not specify the intended relation-length condition or operator norm. The supplied open status is dated curation, not an independently verified conclusion.
\subsection{Sources}
\begin{itemize}
\item Local UnsolvedMath registry, supplied complete source record: permanent ID \(30004338\), exact source problem code \texttt{catalogue entry}, title ``Free Quasi-Actions and Almost-Ramanujan Subgroups,'' in the Oberwolfach Reports Open Problems collection. No individual proposer is supplied; the proposed year is recorded as 2019.
\item Supplied report attribution: \emph{Groups, Dynamics, and Approximation} (2020), Oberwolfach Reports, DOI 10.4171/owr/2019/55, pp.\ 3409--3458. \url{https://doi.org/10.4171/owr/2019/55}. The supplied background describes these reports as documenting mathematical workshops and their open-problem discussions. No problem-specific page or passage is supplied.
\item Supplied literature source, ``OWR report on groups, dynamics, and approximation'': \url{https://oa.tib.eu/renate/server/api/core/bitstreams/f27753e6-ce45-4d59-ba59-812d45c5f07b/content}.
\item Supplied literature source, ``Alon--Wei, large-girth almost-Ramanujan graphs'': \url{https://arxiv.org/abs/2304.01281}. This label and attribution are reproduced as supplied; their correspondence to the linked work and its relevance have not been independently verified.
\item The exact supplied literature-check date is \texttt{2026-08-22}. The supplied dated curation labels the problem open and assesses both the quasi-action classification and the large-girth free-group Cayley-graph question as unresolved, stating that partial graph constructions settle neither formulation. This is supplied dated curation, not an independently verified literature assessment. The record itself requests independent verification of time-sensitive status and cited claims.
\item The supplied statement status is \texttt{corrected\_verified}; its accompanying qualification says that only formula delimiters, notation, grammar, and typography were repaired, with mathematical content and scope preserved. This label does not establish mathematical correctness or resolve the formulation issues above.
\end{itemize}
\end{document}
